\section{Supplementary figure}
\label{sec:appendix-figures}

 Figure~\ref{fig:model-skill-uk} is the UK/Europe counterpart of Figure~\ref{fig:model-skill}.

\begin{figure}[htbp]
\centering
\includegraphics[width=\textwidth]{model_comparison_rmse_uk_europe}
\caption{As Figure~\ref{fig:model-skill}, for the UK/Europe region:
fp32 forecast skill of the models studied, mean RMSE against ERA5 vs.\
lead time.}
\label{fig:model-skill-uk}
\end{figure}

\section{bf16 compared with fp16}
\label{sec:appendix-bf16}

Figure~\ref{fig:rmse-grid-naive} shows the RMSE of Aurora's bf16 forecasts
when the whole model is converted, so that it also holds its state in bf16
in physical units (\S\ref{sec:results-state}), compared with the bf16
network with the state in fp32. The whole-model conversion raises the
RMSE by 19--167\%, the bf16 network alone by at most 25\%.

\begin{figure}[tbp]
\centering
\includegraphics[width=\textwidth]{aurora_bf16_whole_global}
\caption{Mean global RMSE against ERA5 of Aurora's fp32 forecasts, of its
forecasts with the network in bf16 and the state in fp32, and of its
forecasts with the whole model converted with
\texttt{model.to(torch.bfloat16)}, which also rounds the state in physical
units; 52 dates in 2024.}
\label{fig:rmse-grid-naive}
\end{figure}

\begin{table}[htbp]
\centering
\scriptsize
\setlength{\tabcolsep}{3pt}
\caption{Change in global RMSE (\%) relative to fp32 with the network in bf16 and in fp16 (state in fp32 in both), 52 dates in 2024. Bold: significant at the 5\% level (paired $t$-test).}
\label{tab:bf16}
\begin{tabular}{llrrrrrrrrrrrr}
\toprule
 & & \multicolumn{4}{c}{z500} & \multicolumn{4}{c}{2t} & \multicolumn{4}{c}{msl} \\
Model & Precision & d1 & d3 & d5 & d10 & d1 & d3 & d5 & d10 & d1 & d3 & d5 & d10 \\
\midrule
FastNet & bf16 & \textbf{+1.1} & \textbf{+2.1} & \textbf{+3.0} & -0.6 & \textbf{+0.2} & \textbf{+0.6} & \textbf{+0.8} & -0.2 & \textbf{+1.0} & \textbf{+1.7} & \textbf{+1.9} & \textbf{-2.0} \\
 & fp16 & \textbf{-0.3} & \textbf{-0.3} & +0.1 & 0.0 & \textbf{-0.1} & \textbf{-0.2} & \textbf{-0.2} & -0.1 & \textbf{-0.1} & \textbf{-0.2} & +0.1 & +0.3 \\
\midrule
Pangu-Weather & bf16 & \textbf{+3.1} & \textbf{+1.0} & \textbf{+0.6} & +0.7 & \textbf{+0.7} & \textbf{+0.3} & 0.0 & -0.6 & \textbf{+0.6} & \textbf{+0.5} & +0.4 & +0.4 \\
 & fp16 & \textbf{+0.1} & \textbf{+0.1} & \textbf{+0.1} & \textbf{+0.2} & \textbf{0.0} & \textbf{+0.1} & \textbf{+0.1} & \textbf{+0.2} & \textbf{0.0} & \textbf{+0.1} & \textbf{+0.1} & \textbf{+0.2} \\
\midrule
FourCastNet & bf16 & \textbf{+66.4} & \textbf{+28.9} & \textbf{+16.8} & \textbf{+28.5} & \textbf{+14.6} & \textbf{+19.4} & \textbf{+20.0} & \textbf{+271.4} & \textbf{+11.0} & \textbf{+10.2} & \textbf{+8.5} & \textbf{+23.0} \\
 & fp16 & \textbf{+1.4} & \textbf{+0.6} & \textbf{+0.2} & -0.2 & \textbf{+0.2} & \textbf{+0.4} & \textbf{+0.5} & -0.6 & 0.0 & \textbf{-0.2} & \textbf{-0.3} & \textbf{-0.7} \\
\midrule
SFNO & bf16 & \textbf{+0.9} & 0.0 & 0.0 & +0.1 & \textbf{+0.8} & \textbf{+1.0} & \textbf{+1.0} & +0.4 & \textbf{+2.0} & \textbf{+1.0} & \textbf{+0.4} & -0.1 \\
 & fp16 & \textbf{0.0} & \textbf{0.0} & \textbf{+0.1} & 0.0 & \textbf{0.0} & \textbf{0.0} & \textbf{0.0} & 0.0 & \textbf{0.0} & 0.0 & \textbf{0.0} & 0.0 \\
\midrule
Aurora & bf16 & \textbf{+5.4} & \textbf{+7.2} & \textbf{+7.3} & \textbf{+7.0} & \textbf{+7.7} & \textbf{+21.3} & \textbf{+25.2} & \textbf{+18.3} & \textbf{+6.3} & \textbf{+8.5} & \textbf{+8.1} & \textbf{+7.4} \\
 & fp16 & \textbf{+0.2} & \textbf{+0.1} & \textbf{+0.1} & \textbf{+0.3} & \textbf{+0.1} & \textbf{+0.1} & \textbf{+0.2} & +0.1 & \textbf{+0.1} & \textbf{+0.3} & \textbf{+0.2} & \textbf{+0.5} \\
\bottomrule
\end{tabular}
\end{table}


Table~\ref{tab:bf16} compares the change in global RMSE with the
network in bf16 and in fp16, the state in fp32 in both, on the 52 dates.
For every model, the bf16 penalty exceeds the fp16 change at almost
every lead time, by one to two orders of magnitude for Aurora and
FourCastNet. FourCastNet's bf16 forecasts degrade widely by day 10: its
2t error is more than twice that of fp32 on 21 of the 52 dates. Where
bf16 does better than fp32, the differences are small and, except for
FastNet's msl at day 10 ($-2.0\%$), not significant.

\section{Minimum detectable effect}
\label{sec:appendix-mde}

Let $d_i$ be the fp16 RMSE minus the fp32 RMSE on date $i$, with mean
$\bar d$ and standard deviation $s$ over the $N$ dates, and let
$t_{p,\nu}$ be the $p$-quantile of the $t$ distribution with $\nu$ degrees
of freedom. The two-sided paired $t$-test at the 5\% level finds a
significant change when $|\bar d|$ exceeds $t_{0.975,N-1}\,s/\sqrt{N}$.
If the true change is $\delta>0$, $\bar d$ scatters around $\delta$ with
standard error $s/\sqrt{N}$, and lies above
$\delta-t_{0.8,N-1}\,s/\sqrt{N}$ with probability 0.8. The test therefore
detects the change with probability 0.8 when this value equals the
significance threshold, which gives the minimum detectable effect
\begin{equation}
\delta=\left(t_{0.975,N-1}+t_{0.8,N-1}\right)\frac{s}{\sqrt{N}}.
\label{eq:mde}
\end{equation}
The first term is the significance threshold. A true change of that size
would be detected only about half the time, and the second term is the
margin that raises this to 80\%. The same holds for $\delta<0$ by symmetry.
For $N=52$, $t_{0.975,51}=2.008$ and $t_{0.8,51}=0.849$, so that
$\delta\approx0.40\,s$; Table~\ref{tab:power} gives $\delta$ as a percentage
of the fp32 RMSE.

Equation~\eqref{eq:mde} treats $\bar d$ as following a $t$ distribution
shifted by $\delta$, whereas the test statistic strictly follows a
noncentral $t$ distribution. For $N=52$ the exact calculation gives the
same value of $\delta$ to within 0.1\%. Like the test itself, the
calculation assumes that the $d_i$ are independent between dates, and it
treats the observed $s$ as the true standard deviation.
